\begin{frame}{$\varepsilon=0$의 대응물 — 왜 $a=\mu$인가}
  \begin{columns}[T]
    \begin{column}{0.55\textwidth}
      \textbf{대응 원리}
      \begin{itemize}
        \item \textbf{학습 도중(샘플) 평가}: 학습 중과 같은 $a=\mu+\sigma z$로
          실행 — \textbf{탐험 노이즈가 섞인} 성과.
        \item \textbf{학습 후(결정적) 평가}: $a=\mu_\theta(s)$만 내보내고
          샘플링 안 함 — \textbf{연속제어의 ``greedy 평가'' 대응물}.
      \end{itemize}
      \vskip0.4em
      \kb{왜 $a=\mu$가 $\varepsilon=0$인가?}
      \begin{itemize}
        \item $\varepsilon$-greedy에서 $\varepsilon$의 역할 = ``확률
          $\varepsilon$로 아무 행동이나''.
        \item 정규분포 정책에서 ``아무 행동''에 해당하는 것은 \textbf{샘플링
          노이즈 $z$} — 빼면 ``정책의 최선 추측'' $\mu$만 남는다.
      \end{itemize}
    \end{column}
    \begin{column}{0.42\textwidth}
      \begin{block}{$\sigma$를 0으로 만들면 안 되는 이유}
        $\sigma$는 \textbf{학습된 파라미터} — 11.2의 엔트로피 보너스가
        0으로 쪼그라들지 않게 브레이크 역할.\\
        평가 때 $\sigma=0$으로 ``학습''하는 것이 아니라,
        \textbf{파라미터를 그대로 두고 샘플링만 하지 않는 것}이 정직한
        ``탐험 제거''.
        {\footnotesize (``$\sigma$를 0으로 설정해서 돌려라''는 오답 —
        파라미터를 바꾼 \emph{다른} 정책의 실행이 되기 때문.)}
      \end{block}
    \end{column}
  \end{columns}
\end{frame}
